\section*{Chapitre \ref{ch:m2:thinking-in-expected-value} --- Penser en espérance mathématique}

\begin{solution}{exo:m2:thinking-in-expected-value:1}
L'edge de $0.4 \times 2 - 0.6 = 0.2$ par unité misée~; fraction de Kelly $0.2/2 = 0.1$.
\end{solution}

\begin{solution}{exo:m2:thinking-in-expected-value:2}
L'espérance est dominée par quelques trajectoires extrêmement chanceuses~: la richesse espérée croît comme $1.04^{300}$
alors que la richesse de presque chaque joueur croît comme $e^{300g}$, ce qui est bien plus faible. La médiane
ou le \omterm{def:m2:thinking-in-expected-value:growth}{taux de croissance}, décrit ce qu'obtient un joueur typique.
\end{solution}

\begin{solution}{exo:m2:thinking-in-expected-value:3}
$2c - c^2$~: 75~\% à demi-Kelly, 43.75~\% au quart.
\end{solution}

\begin{solution}{exo:m2:thinking-in-expected-value:4}
$0.2^{2/1 - 1} = 20\%$ à Kelly plein~; $0.2^{3} = 0.8\%$ à demi-Kelly.
\end{solution}

\begin{solution}{exo:m2:thinking-in-expected-value:5}
$c = 2/(1 + \ln 0.1/\ln 0.8) = 0.18$~: moins d'un cinquième de Kelly.
\end{solution}

\begin{solution}{exo:m2:thinking-in-expected-value:6}
$t^2/(t^2 + 1) = 4/5 = 0.8$.
\end{solution}

\begin{solution}{exo:m2:thinking-in-expected-value:7}
Environ 94~\% à demi-Kelly et à Kelly (94.2~\% et 93.7~\% sur 4\,000 joueurs simulés à graine fixée), proche de
la valeur de 95~\% que citent les auteurs pour des fractions constantes de 10 à 20~\%.
\end{solution}

\begin{solution}{exo:m2:thinking-in-expected-value:8}
Avec un ratio de Sharpe de 0.1 sur 200 transactions, l'edge a une $t$-statistique d'environ 1.4~: il est
incertain, et miser la fraction de Kelly estimée en totalité conduit à sur-miser dès que l'estimation est
élevée. La croissance espérée est maximale pour environ les deux tiers de la fraction estimée,
et le Kelly plein comporte aussi une probabilité de 50~\% de voir un jour le capital réduit de moitié.
\end{solution}

\begin{solution}{pb:m2:thinking-in-expected-value:1}
\textbf{1.} Edge de 0.1 par unité misée, fraction de Kelly 0.1, croissance de 0.501~\% par transaction.
\textbf{2.} Moyenne 0.1, écart-type 0.995, ratio de Sharpe 0.1005 par transaction.
\textbf{3.} 50~\%.
\textbf{4.} 0.46 de Kelly.
\textbf{5.} Au deuxième ordre, le \omterm{def:m2:thinking-in-expected-value:growth}{taux de croissance} à deux fois Kelly est nul~: le gain issu de l'edge
est exactement compensé par le frein que la variance impose aux mises plus grandes.
\textbf{6.} La moyenne est estimée avec une erreur type de 0.0704~; $t = 1.42$.
\textbf{7.} $c^* = 0.669$.
\textbf{8.} 0.7, le meilleur des multiples essayés, par pas de 0.1.
\textbf{9.} 19~\%~: 29.1 contre 35.9 points de base par transaction.
\textbf{10.} $0.5^{2/0.669 - 1} = 25\%$.
\textbf{11.} Avec 50 transactions, 0.34~; avec 2\,000, 0.95.
\textbf{12.} La croissance décroît de façon quadratique autour de la fraction de Kelly, mais le frein croît avec
le carré de la mise, et un edge surestimé fait augmenter la mise~; les pertes
dues à la sur-mise comportent une vraie probabilité de ruine, celles de la sous-mise ne font que ralentir la croissance.
\textbf{13.} Il faut la réduire~: des pertes plus fortes que ne le laisse supposer la variance rendent les grosses mises plus coûteuses
que ne le montre l'approximation quadratique.
\textbf{14.} Prendre le plus petit entre le multiple de Kelly après shrinkage et le multiple qui maintient le
\omterm{def:m2:thinking-in-expected-value:frac}{risque de ruine} dans la tolérance fixée.
\textbf{15.} Que les joueurs disposant d'un edge clair, lorsqu'ils peuvent choisir leur mise, sur-misent et jouent de façon erratique~:
une règle de dimensionnement arrêtée à l'avance protège l'edge du trader.
\textbf{16.} Parce que les edges sont estimés avec erreur et s'érodent, les rendements ont des queues épaisses,
les stratégies sont corrélées, et les investisseurs et les contreparties imposent des limites de \omterm{def:m2:thinking-in-expected-value:frac}{drawdown} bien
avant la ruine.
\textbf{17.} Mise à jour bayésienne de l'estimation de l'edge et de son incertitude après chaque
transaction, recalcul de la fraction après shrinkage, avec des limites sur la vitesse à laquelle elle peut changer.
\textbf{18.} Lorsque l'objectif n'est pas la croissance à long terme d'un seul capital de jeu~: un horizon fixé, un cap
sur les gains, une limite stricte de \omterm{def:m2:thinking-in-expected-value:frac}{drawdown}, ou une utilité plus averse au risque que le logarithme.
\textbf{19.} Résultat intitulé~: \emph{l'edge incertain}~: avec l'edge estimé à partir de 200 transactions
à un ratio de Sharpe de 0.1, la meilleure mise vaut les deux tiers de la fraction de Kelly estimée
(0.669 par la formule, 0.7 en simulation), ce qui rapporte 19~\% de croissance de plus que le Kelly plein.
\textbf{20.} Parce que l'edge est une estimation, et que la sur-mise en coûte plus que
la sous-mise.
\end{solution}

\begin{solution}{iq:m2:thinking-in-expected-value:1}
Vingt pour cent~: la fraction de Kelly $p - q/b = 0.6 - 0.4$, qui maximise à long terme le
\omterm{def:m2:thinking-in-expected-value:growth}{taux de croissance}, soit environ 2~\% par lancer. Miser davantage fait croître le frein dû à la variance plus vite que le
gain~; au-delà d'environ 39~\% de la richesse misée, celle-ci diminue à long terme.
\iqlookfor{edge par rapport à la \omterm{def:m2:thinking-in-expected-value:edge}{cote}, et pourquoi ne pas miser davantage.}
\end{solution}

\begin{solution}{iq:m2:thinking-in-expected-value:2}
L'espérance de la richesse pousse à tout miser, puisque quelques issues énormes la dominent~;
l'espérance du logarithme de la richesse est le \omterm{def:m2:thinking-in-expected-value:growth}{taux de croissance} qu'atteint presque chaque trajectoire sur de nombreux paris,
et elle pénalise la variance. La maximiser donne Kelly~; maximiser la moyenne mène à la ruine.
\iqlookfor{issue typique contre issue moyenne.}
\end{solution}

\begin{solution}{iq:m2:thinking-in-expected-value:3}
Sur une courte période $dt$, le rendement logarithmique de la richesse, avec une fraction $f$ dans l'actif (et
le reste au taux zéro) a pour espérance $f\mu\,dt - \tfrac12 f^2\sigma^2 dt$~; la maximisation donne $f^* =
\mu/\sigma^2$ et une croissance $\mu^2/(2\sigma^2)$, soit la moitié du carré du ratio de Sharpe.
\iqlookfor{correction d'Itô et terme quadratique.}
\end{solution}

\begin{solution}{iq:m2:thinking-in-expected-value:4}
Le demi-Kelly conserve les trois quarts de la croissance avec la moitié de la volatilité~; il protège contre
les edges surestimés et les queues épaisses~; il maintient les \omterm{def:m2:thinking-in-expected-value:frac}{drawdowns} dans ce que tolèrent les investisseurs et
les limites de risque~: la probabilité de voir un jour le capital réduit de moitié tombe de 50~\% à 12.5~\%.
\iqlookfor{croissance contre risque, et erreur d'estimation.}
\end{solution}

\begin{solution}{iq:m2:thinking-in-expected-value:5}
Les poids de Kelly valent $\mu_i/\sigma_i^2 = \mathrm{SR}_i/\sigma_i$~: à volatilités égales,
deux fois plus de capital pour la stratégie de ratio de Sharpe 2~; en général proportionnels au ratio de Sharpe sur
la volatilité, puis réduits pour tenir compte de l'erreur d'estimation et des limites.
\iqlookfor{$\Sigma^{-1}\mu$ sans corrélation.}
\end{solution}

\begin{solution}{iq:m2:thinking-in-expected-value:6}
Chaque jour, estimer l'edge et le risque de chaque stratégie à partir de son historique en production et en backtest, avec
leur incertitude~; calculer la fraction de Kelly après shrinkage, la plafonner par la règle de \omterm{def:m2:thinking-in-expected-value:frac}{drawdown}, par
les corrélations avec les autres stratégies et par la capacité~; lisser les variations~; publier les
allocations avec leurs données d'entrée~; passer outre en cas de dépassement~; et journaliser tout pour revue.
\iqlookfor{estimation, shrinkage, limites, lissage, audit.}
\end{solution}
